\chapter
[\CO{[PSI]} Psychologie und psychische Betreuung]
{\CC{[PSI]} Psychologie und psychische Betreuung}
\label{chp:PSY}

\ChapterIntro
{%
~~~
}
{%
\begin{description}
  \item[Maintainer:] Sebastian Gabriel
  \item[Autoren:] Michael Auer, Sebastian Gabriel
  \item[Reviewer:] Michael Auer, Standard-Reviewprozess
  \item[Version:] {\DocVersion} ({\DocVersionInfo})
  \item[SHA1:] {\DocGitCommit}
\end{description}

{\TxTeilkapitelAASS}
}



\Layout{2}{Einleitung}
{
~

{\TakeNotesLarge}
}{}{}{}{}

\subsection{Anforderungen und Stress}

\Layout{2}{Belastung, Anforderung, Beanspruchung}
{
\subparagraph{Stichworte} 
(- Kurze Einführungsphase, )  Parallelen
(Beanspruchungsdauer, -intensität) bzw. Unterschied zw.
Mechanischer und psychischer Beanspruchung = interpersonelle
Faktoren (Wahrnehmung, Coping, Expertise) vs. strukturelle
Gegebenheiten (s.u.). - Mehrdeutige Situationen erhöhen
erlebte Beanspruchung - Ruhe- und Alarmphasen, - Frustrane
Erlebnisse (Vorwürfe von Patienten/Angehörige), - Erlebtes
arbeitet emotional und kognitiv weiter, - persönlich
adäquate Balance zwischen professioneller verständnisvoller
Haltung, strukturellen bzw. fachlichen Notwendigkeiten und
persönlichen Grenzen u. Bedürfnissen.

Beanspruchung wird in der Technik laut Definition des
vielseits beliebten Online-Nachschlagewerks Wikipedia als
\Speech{"`die Auswirkung einer äußeren Belastung auf das Innere eines
Körpers bzw. eines Werkstoffes"'} bezeichnet. Mechanische
Beanspruchung ist durch die Ausbildung von mechanischen
Spannungen gekennzeichnet. Vice versa kann psychische
Beanspruchung als die Entstehung psychischer
Spannungszustände verstanden werden, die im größeren Ausmaß
als psychische Belastung bzw. Stresserleben bezeichnet
werden. Im technischen Kontext haben Qualitäten wie
Beanspruchungsdauer (kurz -- lang) bzw.  -intensität (stark --
schwach) einen wesentlichen Einfluss auf das
Beanspruchungsniveau. Auf psychischer Ebene spielen jedoch --
im deutlichen Gegensatz zur Technik -- mehrere vermittelnde
Faktoren eine wesentliche Rolle, nämlich: persönliche
Wahrnehmung,  kognitive Bewertungen und
Verarbeitungsstrategien/Abwehrmechanismen sowie Routine,
Vorerfahrung und Wissen. Diese können in
Beanspruchungssituationen bewusst oder unbewusst zum Tragen
kommen und sollen im Folgenden (s. Stresstheorie) ausgeführt
werden.

Inhaltliche und strukturelle Aspekte der Gesundheitsberufe
stellen diesbzgl. wesentliche Einflussfaktoren dar: zum
Beispiel sieht sich das Fachpersonal häufig mit  komplexen,
mehrdeutigen Situationen konfrontiert, wie z.B. bei
Notfällen und Erkrankungen, die die Handlungsplanung bei der
Versorgung erschweren und damit nicht zuletzt auch die
emotionale Anspannung erhöhen können. Weiters wechseln
während des Dienstes oft Ruhe- bzw. Bereitschaftsphasen, die
der Entspannung dienen können, mit plötzlichen Aktivitäts-
und Alarmphasen, in denen es zu zwangsläufig zu
Stresssituationen kommt. Langfristig kann dieser
unvorhersehbare bzw. unkontrollierbare Wechsel im
Arbeitsrhythmus als unangenehm wahrgenommen werden und zum
Stresserleben beitragen.  Nicht zuletzt stehen Mitarbeiter
helfender bzw. versorgender Systeme nicht nur körperlich,
sondern auch psychisch verletzten Menschen gegenüber. Dies
hat oft zur Folge, dass  das Personal mit leidenden
Äußerungen, aber  auch mit verbalen Vorwürfen und Kritik
durch oft überforderte Betroffene oder Angehörige
konfrontiert wird.

Eine Besonderheit psychischer Beanspruchung stellt die
Tatsache dar, dass emotionale bzw. gedankliche Verarbeitung
stressreicher Situationen bewusst oder unbewusst nachwirkt.
Wiederaufkeimende Gedanken und Gefühle an vergangene
Einsätze halten sich nicht an Dienstzeiten, sodass auch nach
Dienstschluss innerlich emotional ge- und verarbeitet wird.
Die individuell erlebte Beanspruchung hängt auch von der
Fähigkeit des Einzelnen ab, in seiner Rolle als
Professionist im Versorgungssystem eine persönlich
angemessene Balance zwischen einer professionellen
verständnisvollen Haltung (z.\,B. soft skills wie empathisches
Umgehen etc.), strukturellen bzw. fachlichen Erfordernissen
seiner Rollenfunktion als Sanitäter (z.\,B.
Dienstvorschriften, Behandlungsstandards) und persönlichen
Grenzen u. Bedürfnissen (z.B. eigenes Wohlbefinden, Gefühle)
zu finden (Beispiel siehe nächstes Kapitel Rollenkonflikt).
}
{}
{}
{}
{}


\Layout{2}{Rollenkonflikt}
{~
\begin{center}
\FormatTable
\begin{tabularx}{\linewidth}{XcX}
\toprule\addlinespace
Schnelles, aber professionelles und möglichst fehlerfreies
Handeln bei gleichzeitig hoher Verantwortung
&
vs.
&
Emotional berührende Situationen, 

Konfrontation mit den
Themen Unfall/Krankheit, Leid, Trauer, Schock und Tod.
\\\addlinespace\bottomrule
\end{tabularx}
\end{center}



In Gesundheitsberufen Tätige üben nicht nur einen Beruf aus,
sondern übernehmen auch automatisch eine besondere
gesellschaftliche Rolle (Rollenfunktion).  In dieser
Rollenfunktion sehen sie sich mit einem breiten Spektrum an
belastenden Themen konfrontiert, deren erschöpfende
Darstellung hier nicht abgehandelt werden kann. Als Beispiel
soll der Rollenkonflikt, der sich aus der Auseinandersetzung
mit Belastungen im beruflichen Kontext ergibt, erörtert
werden.

\subparagraph{Professionist vs. Person (Verantwortlichkeit vs. persönliche
Gefühle)}

Die Rolle als Angehöriger eines Gesundheitsberufs
beinhaltet, dem Versorgungsauftrag (Versorgung erkrankter
Menschen)  verpflichtet zu sein, d.h. sich für die
Gesundheit bzw. in letzter Instanz evtl. für das Leben eines
anderen Menschen verantwortlich zu fühlen.

Im Rahmen dieser Aufgabe werden dem professionellen Helfer
diverse Handlungs- und Entscheidungskompetenzen
überantwortet – z.B. die Wahl der Maßnahmen am Patienten,
die Organisation des Einsatzes etc. –, zwischen denen er
womöglich auch unter Zeitdruck nach bestem Wissen und
Gewissen entscheiden soll.

Gleichzeitig übt eine fühlende Person diese
verantwortungsvolle Versorgungsrolle aus, die sich mit
berührenden Themen wie Unfall/Krankheit, Leid, Trauer,
Schock und Tod konfrontiert sieht, sodass ihr Verhalten und
Erleben beeinflusst wird durch Gefühle der Anteilnahme,
Belastungs- und Frustrationsgefühle, persönliche
Erinnerungen, Stimmungslagen etc.

Diesen Rollenkonflikt in einer angenehmen und funktionellen
Balance zu halten setzt die Bereitschaft zur \II{Selbstreflexion}
\A{Erklärung!} und Auseinandersetzung mit persönlichen Grenzen
und Emotionen voraus. Diesbzgl. Ungleichgewichte  können
sich einerseits durch das übermäßiges Verhaften in der
Professionistenrolle („Versorgungsroboter“), andererseits
aber auch durch zu hohes Einbringen persönlicher
Gefühlsanteile aufgrund von Gefühlsansteckung (=\,Übernahme
der Gefühle des Betroffenen) bis hin zur Lähmung  durch
emotionalen Kontrollverlust (z.\,B. emotionaler Schock)
ausdrücken.

Eine balancierte Rollenaufteilung „Professionist vs. Person“
stellt eine Ressource (Schutzfaktor) dar, da sie hilft, das
persönliche vom professionellen zu trennen, aber
gleichzeitig auch erlaubt, einen persönlichen Stil in die
Versorgungssituation einzubringen. Dies erfordert  eine
ausgewogene Integration der beiden dargestellten Elemente
„Professionistenrolle“ und „Rolle als emotionales
Individuum“.

Das Bewusstmachen persönlicher Gefühle, Haltungen und
Grenzen -- z.\,B. durch Selbstreflexion, Gespräche etc. --  ist
ist die Voraussetzung, um eine solche Balance für sich
selber zu finden.
}
{}
{}
{}
{}



\label{topic:stress}




\Layout{2}{Stressursachen, -faktoren, -entstehung, Überforderung,
Unterforderung}
{
Physiologie von Stress, evolutionärer Nutzen, Stressmodelle/
-theorien, Aktivierung, individuelles Aktivierungsprofil,
kognitive Bewertung der Gegebenheiten einer Situation,
persönliche Ressourcen, Kritik an „den Stress, Eu-
Distress“, persönliche Ressourcen, Resilienz, Proaktives
Coping, Erkennen und Umgang mit eigenen Schwächen, Zulassen
von Schwächen.

}
{}
{}
{}
{}




\Layout{2}{Stressauswirkungen, Beanspruchungsfolgen, Früherkennung}
{
Psychische/physische Ebene, erhöhte Vulnerabilität,
Reizbarkeit, gedrückte Stimmung, Zynismus, pathologische
Phänomene.
}
{}
{}
{}
{}





\Layout{2}{Grundsätze der Stressvermeidung, Maßnahmen der Verhütung und
Verminderung von Beanspruchungsfolgen}
{
Umgangsregeln im Team, Kommunikationsregeln, Rituale,
interpersonell intrapersonell, formell informell,
Stressverarbeitungsstrategien kognitiv, emotional,
Ablenkung, Kohärenzgefühl durch Kompetenz, Entlastung durch
zwischenmenschliche Ressourcen, intrapersonelle Ressourcen,
soziale Unterstützung.
}
{}
{}
{}
{}

\subsection{Wenn die Anforderungen zum Problem werden:
Psychische Belastungssyndrome}

\subsubsection{PTSS}
\label{topic:ptss}


% \subsection{Belastungen im Dienst}
% In den Teil TAE umgesiedelt [GABS, 2010-08-08]






\subsubsection{Burn-out}
\label{topic:burn-out}

{\TakeNotesLarge}

\opt{STATUS-EXPERIMENTAL}{
Karutz, H.: Wenn die Belastungsgrenze erreicht ist: Psychologische Selbsthilfe
in Extremsituationen Rettungsdienst, 2009, 2009/12, 40-45 \cite{Karutz200912}
}

\subsection{Psychische Betreuung von Kranken/Verletzten}
Gesprächsführung
    \FullPrettyRef{topic:kommunikationsregeln},
    \FullPrettyRef{topic:umgang-psy}

\subsection{verwirrte Patienten}

    wird integrativ in allen Kapitel abgehandelt
\subsection{Begleitung und Betreuung Sterbender}

Dafür gibt es kein Patentrezept und keinen "`goldenen Weg"'.
Wieso? 

Der Tod ist ein integraler BEstandteil des Lebens. So banal
diese Aussage auch scheinen mag, so ambivalent ist das
Verhalten des bzw. der Menschen zu ihm. Nach tausenden
Jahren stellt der Tod und der Sterbeprozess noch immer
(oder vielleicht sogar mehr denn je) eine große
Herausforderung sowohl für den Sterbenden, als auch für
dessen Umgebung dar. 

\expandedtitle{simple}{\Speech{"`Der Tod ist etwas
persönliches."'}}

Die Auseinandersetzung mit dem Tod und er eigenen
Endlichkeit verläuft bei jedem Menschen äußerst
unterschiedlich und ist von dessen Weltbild, Umgebung,
Lebensgeschichte und -situation und vielen anderen Faktoren
abhängig, und kann im Verlauf des Lebens ebenso radikalen
Änderungen unterworfen sein. Mögliche Bewältigungstrategien
sind:

\begin{itemize}
  \item Verdrängung
  \item Hoffnungen, welche sich an den Tod anknüpfen
  (Wiedergeburt, Leben im Jenseits)
  \item Vermächtnis an die Nachwelt; Schaffen von etwas, was
  "`größer als man selbst ist"'
  \item Rationalisierung (Vorsorgeversicherungen,
  Testamente, \ldots)
  \item Rituale (Gebete, Begräbniszeremonien,
  Leichenschmaus, \ldots)
\end{itemize}

\subsubsection{Der todkranke und sterbende Patient}

Rettungsdienst: Was erwartet der Patient und dessen
Angehörigen?

\expandedtitle{roundbox}{\Speech{"`Das große Problem hat das
Fachpersonal."'}}

\subsubsection{Der unerwartete Tod}
\begin{itemize}
  \item Das Problem des Patienten bleibt grundsätzlich das
  Problem des Patienten
  \item Der Helfer muss im Rahmen der Möglichkeiten im Sinne
  des Patienten handeln und seine angemessenen
  Bedürfnisse befriedigen.
  \item Wenn der Helfer mit unzufireden mit der Betreuung
  ist, muss er versuchen, den Rahmen des Möglichen
  zu erweitern.
\end{itemize}

\begin{itemize}
  \item Man darf nicht seine eigenen Vorstellungen
  bezüglich der Hilfe mit denen des Patienten
  gleichstellen. Wenn man dem Patienten nutzen möchte, muss
  man wissen, was \E{er} tatsächlich will.
\end{itemize}


\subsection{Supervision}

    wird extern abgehandelt (PEER-Vortrag)
























\subsection{Psychische Betreuung von kranken/verletzten Menschen}

\Layout{2}{Beziehungsaufbau generell}
{

}
{}
{}
{}
{}

\Layout{2}{Umgang mit Psychiatrisch
erkrankten Menschen}
{

}
{}
{}
{}
{}

\Layout{2}{Suchterkrankte Menschen}
{

}
{}
{}
{}
{}


\subsection{Coping}

{\TakeNotesLarge}

